\section{从大科学到创业公司：为什么现在能做}

前面三章讲物理和工程，本章回答组织问题：为什么核聚变可以由创业公司推进。杨钊的判断是，聚变在科学可行性和不计成本工程可行性上已经有基础；真正缺的是商业可行性，也就是在最短时间、最低成本下把度电成本降到可接受范围。这个阶段更像 SpaceX 式工业品创业：不是发现物理定律，而是把复杂系统从实验室推向低成本工程化。

创业公司适合做这件事，是因为它可以缩短决策链路、自己掌握关键系统、快速迭代、不断把成本压向原材料和充分竞争供应链。能量奇点的选择是核心系统自研，尤其是磁体；这不是因为所有东西都必须自己造，而是因为如果关键子系统是黑盒，就无法知道未来要优化成本和性能时到底该动哪里。

\begin{figure}[H]
\centering
\includegraphics[width=0.94\textwidth]{public-private-fusion.png}
\caption{公共大科学 vs 私营创业：核聚变从国家工程走向创业公司竞速。自制概念图，依据 00:34:00--00:41:03 对谈内容整理。}
\end{figure}

\begin{knowledgebox}{读图：大科学解决可行性，创业公司压缩成本曲线}
公共大科学适合长期基础验证和国际协作，私营创业适合围绕清晰商业目标做工程压缩。两者不是互斥关系：没有公共装置积累，就没有创业路线的物理基础；没有创业公司压成本，聚变也可能长期停留在“能做但太贵”的状态。
\end{knowledgebox}

\subsection{调研 checklist：人才、技术、原料}

21 年创业前，团队调研的关键问题是：在中国推进高温超导托卡马克，会不会被人才、技术和原料卡住脖子。结论是没有明显 no-go：高温超导路线很新，公开知识和学术积累足够起步；中国工程师团队和制造供应链有红利；核心原材料和供应商体系并非完全不可得。真正困难在于大量细碎工程问题要靠自己积累，而不是等外部成熟产业链供应。

\begin{knowledgebox}{创业调研三问}
\begin{tabularx}{\textwidth}{p{0.18\textwidth}p{0.34\textwidth}X}
\toprule
问题 & 调研结论 & 对创业路线的影响 \\
\midrule
人才 & 研究人才稀缺，但中国工程师红利明显 & 团队要把物理设计和工程实现结合起来。 \\
技术 & 高温超导托卡马克是新路线，大家都在摸索 & 先手经验和装置迭代极其重要。 \\
原料 & 许多上游可从原材料和充分竞争供应链获取 & 核心系统自研，非核心环节外部加工。 \\
\bottomrule
\end{tabularx}
\end{knowledgebox}

\begin{figure}[H]
\centering
\includegraphics[width=0.96\textwidth]{fusion-investment-logic.png}
\caption{聚变投资逻辑：巨额资本赌的是文明能源上限，而非短期现金流。自制概念图，依据 00:51:11--00:58:14 与 01:33:00--01:35:00 对谈内容整理。}
\end{figure}

\begin{knowledgebox}{读图：资本在这里买的是技术里程碑}
核聚变创业不是 SaaS 式短周期现金流，而是里程碑驱动：装置建成、第一等离子体、高场磁体验证、Q 大于 10、示范电站、商业复制。每个里程碑都会改变下一轮资金、人才和供应链信心。
\end{knowledgebox}

\subsection{本章小结}

创业公司进入聚变，不是因为物理变简单，而是因为问题从“能否发生”转向“如何低成本商业化”。这个阶段需要快速决策、核心系统自研、供应链打穿和资本耐心。

